\chapter{Reproducibilidad y acceso a los artefactos}
\label{an:reproducibilidad}
% FINAL_HOLDOUT_V2_RESULTS.md §§1,11; TFM_CLOSEOUT.md, entrega de archivo.
\pendiente{Transferir identidades completas, manifests, hashes y comandos
verificados del experimento cerrado. Añadir acceso al archivo de entrega
cuando esté disponible; una ruta local no garantiza acceso al lector.}

\section{Comandos representativos de las interfaces de línea de comandos}
Este anexo permite localizar las cuatro interfaces principales sin reproducir
todos sus argumentos. La referencia operativa completa y mantenida es
\path{docs/USER_GUIDE.md}; la opción \texttt{--help} de cada interfaz muestra
el detalle vigente. Los comandos se ejecutan desde la raíz del repositorio.
Los destinos reales deben ser nuevos cuando así lo exige la guía y una
operación con modelo necesita autorización y credenciales; los ejemplos de
consulta y validación siguientes no ejecutan el experimento de nuevo.

\subsection*{Pipeline CLI}
\textbf{¿Para qué sirve?} Para ejecutar de forma reproducible las etapas
que adquieren y preparan documentos, extraen información, materializan Silver,
construyen o actualizan la referencia territorial y generan Gold mediante
\emph{downstream}.

\begin{lstlisting}[language=bash,title={Ejemplo:}]
uv run python -m renewables_permitting.pipeline --help
\end{lstlisting}
\textbf{Resultado:} ayuda de la Pipeline CLI y lista de operaciones disponibles;
el ejemplo no descarga ni materializa datos.

\subsection*{Admin CLI}
\textbf{¿Para qué sirve?} Para inspeccionar incidencias y registrar o aplicar
decisiones humanas controladas sobre extracciones y posibles antecedentes. No
edita Gold ni las tablas derivadas y no ejecuta por sí sola el pipeline.

\begin{lstlisting}[language=bash,title={Ejemplo:}]
uv run python -m renewables_permitting.admin list \
  --extraction-snapshot <EXTRACCION_VALIDADA>
\end{lstlisting}
\textbf{Resultado:} lista validada de casos pendientes que una persona puede
examinar antes de registrar una decisión.

\subsection*{Evaluation CLI}
\textbf{¿Para qué sirve?} Para validar o congelar la referencia y el evaluador
y ejecutar offline la comparación de las predicciones preservadas frente a la
referencia congelada. No llama al modelo de extracción.

\begin{lstlisting}[language=bash,title={Ejemplo:}]
UV_OFFLINE=1 uv run python -m evaluation.final_holdout_v2.cli \
  evaluate --truth <REFERENCIA_CONGELADA> \
  --predictions <PREDICCIONES_PRESERVADAS> \
  --evaluator <EVALUADOR_CONGELADO> \
  --execution-record <REGISTRO_EJECUCION> \
  --output <NUEVO_DIRECTORIO_RESULTADOS>
\end{lstlisting}
\textbf{Resultado:} un nuevo informe de evaluación offline, o un fallo cerrado
si las entradas y sus identidades no cumplen el protocolo.

\subsection*{Controlador de ejecución primaria}
\textbf{¿Para qué sirve?} Para preparar y controlar la ejecución final del
holdout con estado durable por documento. Registra qué casos terminaron,
fallaron o quedaron en un estado que no debe reintentarse de forma ambigua;
\texttt{TERMINAL\_SUCCESS} indica un artefacto válido, no una predicción correcta.

\begin{lstlisting}[language=bash,title={Ejemplo:}]
env -u PYTHONPATH PYTHONDONTWRITEBYTECODE=1 \
  .venv/bin/python -m evaluation.primary_execution.cli \
  primary-status --run-root <DIRECTORIO_EJECUCION>
\end{lstlisting}
\textbf{Resultado:} resumen de estados terminales, pendientes o indeterminados;
el comando no llama al modelo. El procedimiento completo, incluidas preparación,
autorización y finalización, se encuentra en la guía de usuario y en
\path{docs/evaluation/PRIMARY_EXECUTION_V2.md}.

\chapter{Contrato documental y material de apoyo}
\label{an:contrato}
% architecture/silver_extraction_contract.md; contrato V2; MAPA_FUENTES.md.
\pendiente{Incluir solo el detalle contractual y de anotación necesario para
la revisión científica, un glosario y, si aportan claridad, ejemplos ya
verificados. Evitar trasladar la documentación del código a la memoria.}
